\chapter{Descrizione del progetto}
\label{ch:descrizione-progetto}

\section{Strumenti di sviluppo, versioning e hosting}

Il progetto adotta un approccio rapido e sperimentale, supportato da strumenti AI per refactoring e generazione di codice sotto supervisione umana, con codice versionato su GitHub e architettura containerizzata (Docker) per la portabilità.
L'hosting è valutato tra soluzioni PaaS e VPS; per il training si prevede il ricorso a GPU cloud (es.\ RunPod).
Nella pratica il prototipo attuale è interamente in Python (Flask + PostgreSQL, in produzione su Heroku), come documentato nel Capitolo~\ref{ch:stato-implementazione}; una demo pubblica è disponibile all'indirizzo \url{https://dollarpunk-26786f215068.herokuapp.com/dashboard}.

\section{Raccolta dati}

La previsione dei movimenti di mercato richiede grandi volumi di dati testuali da fonti eterogenee: testate finanziarie (Bloomberg, Reuters, WSJ, Yahoo Finance, ANSA, Il Sole 24 Ore), social media e community (Twitter, Reddit---con relative metriche di engagement---, Telegram, YouTube) e motori di ricerca (SerpApi, \texttt{ddgs}), su temi che spaziano dal sentiment finanziario ai trend macroeconomici, earning report, ESG e criptovalute.
L'acquisizione avverrà via API quando possibile e via scraping (es.\ libreria \texttt{newspaper}) altrimenti, con una strategia comune di strutturazione dei dati.

\section{Analisi dati testuali}

Il nucleo della pipeline è una funzione che elabori un testo e restituisca un oggetto strutturato con i ticker menzionati e, per ciascuno, un \emph{relevance score} e un \emph{sentiment score}:

\begin{lstlisting}[language=json]
{
"title": "If You Invested $1,000 In Lululemon Stock In 2013...",
"source": "Motley Fool",
"time_published": "2023-04-30T13:34:00",
"stocks_mentioned": ["LULU", "NKE"],
"ticker_sentiments": [
  {"ticker": "NKE", "relevance_score": 0.27,
   "sentiment_score": 0.22, "sentiment": "Somewhat-Bullish"}
]
}
\end{lstlisting}

Di seguito le soluzioni valutate per il calcolo affidabile dei tre campi, in un settore in cui strumenti più efficaci possono emergere anche nel breve termine.

\subsection{Tickers}

L'associazione testo--ticker combina più segnali: contesto della raccolta (query mirate), Named Entity Recognition, pattern matching (\verb|$AAPL|, \verb|(NASDAQ:AAPL)|) e dizionari nome--ticker.
Un sistema robusto integra le fonti con un \emph{confidence score}---alto per citazioni dirette, medio per nomi aziendali, basso per riferimenti indiretti---riducendo l'ambiguità dell'associazione.

\subsection{Relevance score}

Il \textit{relevance score} misura quanto un ticker sia centrale nel testo.
Il conteggio normalizzato delle occorrenze basta per un prototipo; approcci più contestuali includono i pesi di attenzione dei Transformer, i sistemi commerciali che valutano frequenza, posizione e ruolo semantico \parencite{ravenpack_relevance}, o l'assegnazione diretta via prompting di un LLM, resa credibile dai risultati di \textcite{lopezlira2024chatgptforecaststockprice} sullo scoring di headline.

\subsection{Sentiment score}

Il \textit{sentiment score} mappa il tono verso un ticker su una scala continua (da $-1$ a $+1$), tramite modelli specializzati come FinBERT o LLM.
Per maggiore precisione si adotta il sentiment a livello di entità \parencite{tang-etal-2023-finentity}, come nei sistemi commerciali (Refinitiv, Bloomberg) e open-source \parencite{yang2023fingptopensourcefinanciallarge}; in alternativa, LLM zero-shot generano punteggi continui via prompt senza training specifico \parencite{bdcc8110143}.

\subsection{Fine-tuning e RAG}

Il fine-tuning di un modello open-weight consente di ottenere un sistema che restituisca direttamente l'oggetto strutturato (ticker, relevance, sentiment), unificando la pipeline: candidati naturali sono FinBERT \parencite{araci2019finbert} e gli LLM open-weight recenti (LLaMA~3, Qwen), affinabili su dataset con etichette a livello di entità---da challenge pubbliche \parencite{maia201818} o aggregati tramite API di notizie---su GPU cloud.
L'alternativa senza addestramento è RAG \parencite{lewis2020retrieval}: si indicizza il corpus in una base vettoriale e un LLM genera lo scoring dal contesto recuperato, con il vantaggio dell'aggiornabilità immediata dell'indice.

\subsection{Sintesi}

La scelta operativa dipende da due risorse: il dataset etichettato e la capacità di calcolo.
Con un dataset a livello di ticker disponibile---condizione oggi soddisfatta dall'archivio descritto nel Capitolo~\ref{ch:stato-implementazione}---la via maestra è il fine-tuning supervisionato; in alternativa il RAG evita l'addestramento.
In assenza di dataset restano lo zero-shot e l'annotazione manuale o LLM-assistita di un corpus iniziale.

\section{Decay temporale (solo sul relevance)}

Calcolati \textit{tickers}, \textit{sentiment} e \textit{relevance}, i punteggi vanno distribuiti nel tempo per ottenere una misura giornaliera del sentiment attivo.
Si applica un decadimento esponenziale al solo \textit{relevance score}---il sentiment resta invariato, ma pesa in proporzione alla rilevanza residua:
\[
R_{i,d} = R_{i,0} \cdot e^{-\lambda \Delta t}, \qquad S_{i,d} = S_{i,0},
\qquad
\text{AvgW}_d = \frac{\sum_i S_{i,0} \cdot R_{i,d}}{\sum_i R_{i,d}}
\]
dove $\Delta t$ sono i giorni dalla pubblicazione e $\lambda$ un parametro di decay da stimare.
L'approccio è coerente con i provider commerciali, che applicano decay esponenziali alla rilevanza \parencite{ravenpack_relevance, ravenpack_forex, ravenpack_credit}.
Nell'implementazione corrente (Capitolo~\ref{ch:stato-implementazione}) il decadimento è parametrizzato tramite \emph{emivita} $h$, con peso $w = 0{,}5^{\Delta t / h}$: equivalente a $\lambda = \ln(2)/h$, ma direttamente interpretabile.

\subsubsection*{Esempio}

Nella tabella, il relevance viene attenuato giorno per giorno con $\lambda = 0{,}5$ (la news~2 è pubblicata al giorno~1), mentre il sentiment resta fisso:

\begin{center}
\renewcommand{\arraystretch}{1.4}
\begin{tabular}{|c|c|c|c|}
\hline
\textbf{News \#} & \textbf{Giorno 0} & \textbf{Giorno 1} & \textbf{Giorno 2} \\
\hline
1 & S=1.000,\ R=0.800 & S=1.000,\ R=0.485 & S=1.000,\ R=0.294 \\
2 & ---               & S=0.900,\ R=0.700 & S=0.900,\ R=0.425 \\
3 & S=0.750,\ R=0.600 & S=0.750,\ R=0.364 & S=0.750,\ R=0.221 \\
\hline
\textbf{Media ponderata $\text{AvgW}_d$} & $\approx$0.893 & $\approx$0.896 & $\approx$0.896 \\
\hline
\textbf{Massa attiva $\sum_i R_{i,d}$} & 1.400 & 1.549 & 0.940 \\
\hline
\end{tabular}
\end{center}

L'esempio illustra le proprietà strutturali del modello, che formalizziamo e dimostriamo di seguito.

\begin{proposizione}[Proprietà del modello di decadimento]
\label{prop:decay}
Sia $R_{i,d} = R_{i,0}\, e^{-\lambda \Delta t_i}$ con $\lambda > 0$ e $S_{i,d} = S_{i,0}$. Allora:
\begin{enumerate}
  \item[(i)] \emph{(unicità)} $R(t) = R_0 e^{-\lambda t}$ è l'unica soluzione di $R' = -\lambda R$ con dato iniziale $R_0$, ed è equivalente alla parametrizzazione in emivita $w = 0{,}5^{\Delta t / h}$ con $h = \ln 2/\lambda$;
  \item[(ii)] \emph{(confinamento)} per ogni giorno $d$: $\min_i S_{i,0} \leq \text{AvgW}_d \leq \max_i S_{i,0}$;
  \item[(iii)] \emph{(invarianza)} a insieme di notizie fissato, $\text{AvgW}_d$ è invariante al passare del tempo a meno del riordino relativo dei pesi, mentre la massa attiva decade esattamente: $\sum_i R_{i,d+1} = e^{-\lambda} \sum_i R_{i,d}$.
\end{enumerate}
\end{proposizione}
\begin{proof}
(i) Unicità dalla teoria delle equazioni lineari del primo ordine; $e^{-\lambda t} = 2^{-t/h} \iff \lambda = \ln 2 / h$.
(ii) $\text{AvgW}_d = \sum_i \alpha_i S_{i,0}$ con $\alpha_i = R_{i,d}/\sum_j R_{j,d} \geq 0$ e $\sum_i \alpha_i = 1$: combinazione convessa dei sentiment osservati.
(iii) Tra due giorni consecutivi senza nuove notizie ogni peso è moltiplicato per il medesimo fattore $e^{-\lambda}$, che si semplifica tra numeratore e denominatore di $\text{AvgW}$ lasciandola invariata, mentre la somma dei pesi ne risulta contratta esattamente di $e^{-\lambda}$.
\end{proof}

La conseguenza operativa della Proposizione~\ref{prop:decay} è visibile nella tabella: la media ponderata \emph{non} decade (si sposta solo verso le notizie più recenti quando ne arrivano), mentre a decadere è la massa attiva, che cresce al giorno~1 con la news~2 e si dissipa al giorno~2.
Le due grandezze sono complementari---la media misura la \emph{direzione} del sentiment attivo, la massa la sua \emph{intensità}---e un segnale operativo completo le richiede entrambe.

\subsection{Sviluppi futuri: diffusione semantica}

Il decadimento esponenziale è la soluzione dell'equazione $\frac{dR}{dt} = -\lambda R$, che tratta però ogni notizia come indipendente.
In realtà molte news condividono contesto semantico, e trattarle come isolate ignora gli effetti informativi correlati.
L'estensione proposta introduce un termine di diffusione nello spazio semantico degli embedding:
\[
\frac{\partial R(x, t)}{\partial t} = D \nabla^2 R(x, t) - \lambda R(x, t),
\]
con $D$ coefficiente di diffusione semantica: la rilevanza non solo decade, ma si propaga verso contenuti affini, trasformando il relevance da punteggio statico a campo dinamico.

La formulazione così enunciata lascia però aperta una questione di buona posizione: se la rilevanza migra dal punto $x_j$ al punto $x_i$ mentre i sentiment restano fissi, la massa in arrivo amplifica il peso del sentiment \emph{locale} $S_i$, anche quando l'informazione che l'ha generata aveva tono opposto.
La risoluzione naturale è far diffondere, con lo stesso operatore, anche la \emph{massa di sentiment} $m = S \cdot R$, definendo il sentiment come rapporto $S(x,t) = m(x,t)/R(x,t)$: la rilevanza che fluisce porta con sé il proprio tono.
La versione discreta del modello, con le garanzie formali di correttezza, è sviluppata nel Capitolo~\ref{ch:direzioni-future}.

Quanto all'innovatività: la maggior parte degli studi tratta le news come unità indipendenti; i lavori più vicini mostrano la diffusione del sentiment su reti informative \parencite{wan2021sentimentdiffusion} o modellano le interazioni tra notizie per la previsione \parencite{wang2024finin}, ma l'uso esplicito di un'equazione di diffusione--decadimento interpretabile risulta poco documentato in ambito NLP-finance.

\section{Analisi esplorativa: sentiment e prezzo}

Dopo raccolta e scoring, va verificato se il sentiment estratto abbia un legame con i prezzi: grafici comparativi sentiment--chiusure, scatter plot con le variazioni percentuali e correlazioni (Pearson, Spearman) anche con ritardi temporali.
Studi come \textcite{cristescu2023wavelet} rilevano correlazioni significative tra sentiment delle notizie e prezzi di titoli come Apple, Microsoft e Tesla.
La fase è iterativa: senza segnali si rivedono dati o scoring; con una relazione anche parziale si procede verso modelli più sofisticati.

\subsection{Backtesting e simulazione di portafoglio}

Il backtesting simula una strategia su dati storici: ad esempio, long se il sentiment è positivo, flat (o short) altrimenti, con ingresso il giorno successivo al segnale.
Le metriche di valutazione sono rendimento cumulato contro benchmark, accuratezza direzionale e Sharpe ratio semplificato.
L'estensione multi-ticker aggiunge capitale iniziale, allocazione, regole di uscita e costi di transazione: l'obiettivo non è ottimizzare, ma verificare che la strategia regga in condizioni credibili prima di investire in modelli predittivi.

\subsection{Integrazione con modelli predittivi}

Il passo successivo è un modello che utilizzi congiuntamente prezzi, volumi, indicatori tecnici e sentiment su finestre temporali mobili.
Le reti LSTM sono lo standard consolidato per la previsione finanziaria \parencite{fischer2018deep}, e diversi studi \parencite{ouf2024lstm, gu2024finbert} mostrano che l'aggiunta del sentiment come variabile esplicativa migliora la previsione rispetto ai soli dati di prezzo, catturando movimenti spiegati dall'informazione testuale.
Le alternative fondazionali più recenti sono discusse nel Capitolo~\ref{ch:direzioni-future}.
